% geometry/part_intro-DE.tex

\chapter*{Geometrie}
\phantomsection
\addcontentsline{toc}{chapter}{Geometrie}

Der aus den Ingenieurbegriffen übernommene Übergang ist keine anonyme Kurve
mehr. Er ist ein vorgeschlagener Zustand eines identifizierten Alignments.
Seine Berlinish-Komposition bleibt prüfbar, seine Ableitung besitzt Provenienz,
und die Realisierung seiner Bewertung ist angegeben. Berechnung,
Anwendbarkeit und Autorität bleiben getrennt. Die Geometrie muss nun eine
engere Frage beantworten:

\begin{quote}
Welche Struktur gehört intrinsisch zu diesem Alignment, und was erhält erst
durch Repräsentation, Platzierung, Realisierung, Beobachtung oder betrieblichen
Zweck Bedeutung?
\end{quote}

\section*{Die beim Alignment verbleibende Struktur}

Elementfolge, Orientierung, intrinsische Längsparametrisierung und
Krümmungsgesetz gehören zur konstruktiven Beschreibung des
Alignment-Zustands. \emph{Stationierung} bezeichnet in diesem Teil diese
intrinsische Parametrisierung durch $s$, gewöhnlich als Bogenlänge.
Betriebliche Kilometrierung und dargestellte Stationsangaben dürfen $s$
adressieren; sie ersetzen es nicht und werden nicht zu Alignment Identity.

Eine Anfangspose und das intrinsische Gesetz rekonstruieren eine lokale
kartesische Plankurve. Diese Kurve ist bereits sinnvolle Ingenieurgeometrie.
Eine Koordinatenrepräsentation kann sie ausdrücken, und eine spätere
geografische Platzierung kann sie auf einen projizierten oder erdbezogenen
Zusammenhang beziehen, ohne das lokale Alignment minderwertig oder fiktiv zu
machen. Es gilt die Grenze
\[
\begin{aligned}
\text{intrinsische Geometrie}&\neq\text{Koordinatenrepräsentation},\\
\text{geografische Platzierung}&\neq\text{physische Realisierung}.
\end{aligned}
\]

Abbildung~\ref{fig:alignment-object-stationing-curvature} verfolgt denselben
Gerade--Übergang--Bogen-Zustand durch diese Lesarten. Objektbezug und geordnete
Stationsbereiche bleiben stabil; Plankurve und Krümmung legen verschiedene
Folgen derselben Konstruktion offen.

\begin{figure}[htbp]
  \centering
  \resizebox{\textwidth}{!}{% figures/geometry/alignment_object_stationing_curvature-DE.tex
\begin{tikzpicture}[
  x=1cm,y=1cm,>=Stealth,font=\scriptsize,
  panel/.style={draw=black!45,rounded corners=2pt,fill=black!2,
    minimum width=4.05cm,minimum height=3.25cm},
  flow/.style={->,semithick,draw=black!60},
  note/.style={align=center,text=black!75}
]
  \node[panel] (object) at (0,0) {};
  \node[font=\scriptsize\bfseries,anchor=north] at ([yshift=-4pt]object.north)
    {identifizierter Alignment-Zustand};
  \draw[very thick] ([xshift=-1.55cm]object.center) -- ([xshift=1.55cm]object.center);
  \fill ([xshift=-1.55cm]object.center) circle (1.4pt);
  \fill ([xshift=-0.45cm]object.center) circle (1.4pt);
  \fill ([xshift=0.75cm]object.center) circle (1.4pt);
  \fill ([xshift=1.55cm]object.center) circle (1.4pt);
  \node[note,anchor=north] at ([yshift=-7pt]object.center)
    {$I_G$ Gerade\quad $I_U$ Übergang\quad $I_B$ Bogen\\geordnet durch intrinsisches $s$};
  \node[note,anchor=south] at ([yshift=4pt]object.south)
    {Identität, Konstruktion, Provenienz};

  \node[panel,right=0.55cm of object] (plan) {};
  \node[font=\scriptsize\bfseries,anchor=north] at ([yshift=-4pt]plan.north)
    {lokale Planrekonstruktion};
  \draw[->,thin] ([xshift=-1.55cm,yshift=-1.05cm]plan.center) -- ++(3.15,0) node[right] {$x$};
  \draw[->,thin] ([xshift=-1.35cm,yshift=-1.25cm]plan.center) -- ++(0,2.25) node[above] {$y$};
  \draw[very thick]
    ([xshift=-1.35cm,yshift=-.45cm]plan.center) -- ([xshift=-.55cm,yshift=-.45cm]plan.center)
    .. controls ([xshift=.15cm,yshift=-.42cm]plan.center) and ([xshift=.40cm,yshift=.35cm]plan.center)
    .. ([xshift=1.30cm,yshift=.65cm]plan.center);
  \node[note,anchor=south] at ([yshift=4pt]plan.south)
    {$\gamma(s)$ aus Anfangspose und $\kappa(s)$};

  \node[panel,right=0.55cm of plan] (band) {};
  \node[font=\scriptsize\bfseries,anchor=north] at ([yshift=-4pt]band.north)
    {Krümmungsband};
  \draw[->,thin] ([xshift=-1.55cm,yshift=-.80cm]band.center) -- ++(3.15,0) node[right] {$s$};
  \draw[->,thin] ([xshift=-1.35cm,yshift=-1.00cm]band.center) -- ++(0,2.05) node[above] {$\kappa$};
  \draw[very thick] ([xshift=-1.35cm,yshift=-.80cm]band.center) -- ([xshift=-.55cm,yshift=-.80cm]band.center)
    -- ([xshift=.55cm,yshift=.55cm]band.center) -- ([xshift=1.35cm,yshift=.55cm]band.center);
  \draw[densely dashed,black!45] ([xshift=-.55cm,yshift=-1.00cm]band.center) -- ([xshift=-.55cm,yshift=.85cm]band.center);
  \draw[densely dashed,black!45] ([xshift=.55cm,yshift=-1.00cm]band.center) -- ([xshift=.55cm,yshift=.85cm]band.center);
  \node[note,anchor=south] at ([yshift=4pt]band.south)
    {$0$\quad veränderlich\quad konstant};

  \node[panel,right=0.55cm of band] (qual) {};
  \node[font=\scriptsize\bfseries,anchor=north] at ([yshift=-4pt]qual.north)
    {angefügte Qualifikationen};
  \node[note] at ([yshift=.55cm]qual.center) {Koordinatenrepräsentation};
  \node[note] at ([yshift=.05cm]qual.center) {geografische Platzierung};
  \node[note] at ([yshift=-.45cm]qual.center) {physische Realisierung};
  \node[note,anchor=south] at ([yshift=4pt]qual.south)
    {Zweck, Beobachtung, Betrieb};

  \draw[flow] (object) -- (plan);
  \draw[flow] (plan) -- (band);
  \draw[flow] (band) -- (qual);
\end{tikzpicture}
}
  \caption{Ein identifizierter Alignment-Zustand als geordnete
  Stationsbereiche, rekonstruierte lokale Plangeometrie und Krümmungsband.
  Repräsentation, geografische Platzierung und physische Realisierung bleiben
  angefügte Qualifikationen statt intrinsischer Identität.}
  \label{fig:alignment-object-stationing-curvature}
\end{figure}

\section*{Warum Krümmung für die Eisenbahn besonders wichtig ist}

Ein Eisenbahnfahrzeug ist durch das Schienenpaar geführt. Anders als ein
Straßenfahrzeug kann es normalerweise nicht innerhalb eines Korridors seitlich
einen anderen Weg wählen, um die Krümmung des Alignments zu umgehen; dieser
Vergleich behauptet nicht, Straßenkorridorentwurf sei geometrisch oder
dynamisch einfach. Bei der Eisenbahn übersetzt Geschwindigkeit räumliche
Änderung in zeitabhängige Fahrzeuganregung. Überhöhung kann einen Teil der
Querwirkung kompensieren, löscht die Krümmung aber nicht. Krümmung, Überhöhung
und ihre Längsraten beeinflussen Fahrzeugreaktion, Komfort,
Rad--Schiene-Kräfte, Gleisbeanspruchung und Verschleiß.

Darum wird das Krümmungsband im nächsten Kapitel zur primären Lesart. Die
Plankurve zeigt, wohin das Alignment führt. Das Band zeigt unmittelbar, wo es
gerade ist, wo Krümmung konstant bleibt, wo sie sich ändert und ob
Elementanschlüsse die erforderliche Stetigkeit bewahren. Die Mathematik beginnt
nicht erneut mit allgemeinen Kurven, sondern legt die intrinsischen Folgen des
bereits identifizierten Übergangszustands offen.
